\section{Task and evidence boundary}
\label{sec:task}
\begin{figure}[t]
\centering
\setlength{\unitlength}{1mm}
\begin{picture}(164,34)
\footnotesize
\put(0,17){\framebox(28,15){\shortstack{Permitted\\current record}}}
\put(34,17){\framebox(28,15){\shortstack{Evidence\\selection}}}
\put(68,17){\framebox(28,15){\shortstack{Shared\\model reader}}}
\put(102,17){\framebox(28,15){\shortstack{Contract and\\support checks}}}
\put(136,17){\framebox(28,15){\shortstack{Mock\\execution}}}
\put(28,24.5){\vector(1,0){6}}
\put(62,24.5){\vector(1,0){6}}
\put(96,24.5){\vector(1,0){6}}
\put(130,24.5){\vector(1,0){6}}
\put(102,0){\framebox(28,10){\shortstack{Full permitted\\record and labels}}}
\put(116,10){\vector(0,1){7}}
\put(0,0){\makebox(96,10)[l]{\shortstack[l]{Same facts, public catalog and output interface.\\Only selection and the stated control vary.}}}
\end{picture}
\caption{The controlled evidence boundary. The time/access view precedes every retrieval adapter. The reader sees selected evidence; evaluation also has the full permitted record and hidden labels. Labels never enter retrieval or the reader. The full-context reference skips bounded selection. A validated positive proposal alone reaches the fictional executor. Correct withholding is scored separately.}
\label{fig:boundary}
\end{figure}

Each query specifies a company world, target account, objective, valid-time cutoff, knowledge-time cutoff, requester permissions and maximum resource units. An event contains its source identifier, world, entity, relation, value, authority, valid time, observed time, required permissions and explicit supersession links. Authority distinguishes observed evidence, approved policy, proposals and hypotheses. Seller strategy and account observations remain logically separate even though the prototype stores both in local files.

The common view admits an event only when it belongs to the query's world, both time cutoffs are satisfied, and all required permissions are present. An eligible later event can supersede an earlier source only for the same world, entity, relation and authority. Independent conflicting current values remain a conflict; a convenient latest value is not silently chosen. This is a deliberately limited event model with valid-from times and explicit replacements, not a complete implementation of arbitrary bitemporal intervals.

A catalog template has an objective, required entity/relation/authority slots, optional literal values, equality constraints, workflow steps and resource units. Four common slots require approved seller role/geography and observed account role/geography, with the two pairs matching. Task-specific slots connect account state or signals to an offering pain and proof. Permission to read a source is distinct from approval of its content. A single non-response is an observation and cannot change approved targeting.

For a concrete contract illustration, development world \texttt{w69bf7897} targets IT managers in the EU. Fictional account Velora-231 matches those attributes and has an observed hiring signal; separate seller records connect its offering and proof to onboarding. With the required permissions and three available units, either the three-step research review or the two-step evidence review is acceptable. A procurement-specific alternative additionally requires a source absent from that view. Removing proof access blocks the hiring proposals, while a qualification-only objective needs only targeting evidence. This is an illustration of the declared fixture rules, not a model result or evidence that hiring predicts a purchase.

Let $V(q)$ denote the complete permitted current view, and $A(q)$ the set of catalog templates whose constraints are satisfied by $V(q)$. Several alternatives may be acceptable; we do not impose an arbitrary unique commercial winner. The reader returns either one template and cited sources, or a request for evidence. A supported recommendation must select a member of $A(q)$, cite at least one current supplied source for every required slot, and include no irrelevant sources. Withholding is correct only when $A(q)$ is empty. Returning nothing on answerable cases is an error.

All methods see the same public catalog and constraint language. The reader first fills a prerequisite-to-source table, with one supplied source or null for every catalog relation, and then selects a template or withholds. Short labels bind exactly to immutable identifiers. The common grammar allows every supplied label in every row, including irrelevant evidence; it does not choose a correct source. Rows required by the selected template become its citations. Missing entries are not filled, and swapped row bindings fail even if their union contains all necessary sources. Rows concerning other templates do not become extraneous citations for the chosen action. Identifier existence and parts of schema validity are enforced by construction, so they cannot be interpreted as empirical gains from memory. Correct source binding, completeness and template applicability remain model decisions. No post-hoc citation repair is performed.

After validation, a successful proposal becomes a action specification with its objective, account, policy versions, sources, allowed next step, mechanism hypothesis and computed resource estimate. Mock research capabilities produce saved receipts from the complete permitted fictional provider view. They cannot revise the model's answer or add citations to make it pass. Each executed step costs one artificial resource unit. These units check workflow consistency; they are not dollars, conversion forecasts or estimated revenue. Mechanisms are catalog-authored hypotheses, not measured causal effects. The reader is stateless between queries: longitudinal cases test decisions over changing records, not autonomous revision of a running action or learned experience reuse.
